\section{INTRODUCTION}
When studying $N$-body celestial systems, selecting an appropriate numerical integration scheme is fundamental to obtaining accurate and meaningful results.
Leapfrog and fourth-order Runge-Kutta (RK4) are routinely employed as standard integrators across modern computational astrophysics, particularly in non-chaotic regimes.

This paper investigates the two-body gravitational problem within the frameworks of Newtonian mechanics and Hamiltonian formalism.
In classical mechanics, the Hamiltonian is a function describing the total energy of a system in terms of its generalised coordinates and conjugate momenta, governing the time evolution of the system through Hamilton's equations.
This formalism is particularly well-suited to the study of orbital dynamics, as it naturally exposes conserved quantities such as energy and angular momentum against which numerical schemes can be rigorously evaluated.

We restrict our analysis to classical mechanics as defined by Newton and Kepler, applying the governing equations of the field accordingly.
A high-performance simulation engine is developed in C++20 to carry out procedurally generated two-body simulations, with the engine exposed to Python 3.12 via a foreign function interface – a common and practical convention in scientific computing workflows.
To make the resulting dual-resolution study tractable, simulations are distributed across a fixed-size worker pool, with each worker independently generating a problem instance and executing all five candidate integrators under matching initial conditions.
Simulation data is then analysed in Python to draw conclusions regarding integrator accuracy and computational efficiency.

Having identified the most suitable standard integrator, we subsequently construct a higher-order scheme and repeat the analysis under equivalent conditions.
Final conclusions are drawn on the basis of energy and angular momentum conservation behaviour, alongside a comparison of computational cost across all tested methods.